% !TeX program = xelatex
% 语言润色版：正文重写为第一人称过程叙述，未改动原 report.tex。
% 提醒：封面若需要保留“本稿由 AI 参与起草”一类说明，请自行加回；当前版本只保留任务书要求的“AI 辅助过程”一章。
\documentclass[UTF8,a4paper,zihao=-4,fontset=none]{ctexart}
\usepackage[left=23mm,right=23mm,top=23mm,bottom=23mm,headheight=15pt]{geometry}
\usepackage{fontspec,graphicx,booktabs,tabularx,longtable,array,amsmath}
\usepackage{float,caption,pdflscape,fancyhdr,enumitem,needspace}
\usepackage[hidelinks,bookmarksnumbered]{hyperref}
\setCJKmainfont[AutoFakeBold=2]{SimSun}
\setCJKsansfont[AutoFakeBold=2]{SimSun}
\setCJKmonofont{SimSun}
\setmainfont{Times New Roman}
\setsansfont{Arial}
\setmonofont{Consolas}
\ctexset{section={format=\Large\bfseries,beforeskip=1.2em,afterskip=.7em},subsection={format=\large\bfseries,beforeskip=1em,afterskip=.5em}}
\linespread{1.22}
\setlength{\parindent}{2em}
\setlength{\parskip}{.25em}
\setlength{\emergencystretch}{2em}
\setlength{\tabcolsep}{5pt}
\renewcommand{\arraystretch}{1.25}
\captionsetup{font=small,labelfont=bf,labelsep=quad,skip=7pt}
\setlist{nosep,leftmargin=2em,topsep=.3em}
\pagestyle{fancy}
\fancyhf{}
\fancyhead[L]{\small 软件工程实验二}
\fancyhead[R]{\small 酒店预订系统：设计与建模}
\fancyfoot[C]{\small\thepage}
\renewcommand{\headrulewidth}{.4pt}
\setcounter{tocdepth}{1}
\newcolumntype{Y}{>{\raggedright\arraybackslash}X}
\newcommand{\pathtext}[1]{\texttt{\detokenize{#1}}}
\newcommand{\ui}[4][78mm]{\begin{figure}[H]\centering\includegraphics[width=.98\linewidth,height=#1,keepaspectratio]{\detokenize{assets_color/#2}}\caption{#3}\label{#4}\end{figure}}
\newcommand{\umlportrait}[4]{\clearpage\begin{figure}[p]\centering\includegraphics[width=163mm,height=225mm,keepaspectratio]{\detokenize{../deliverables/UML_color/#1}}\caption{#2}\label{#3}\par\vspace{2mm}{\small #4}\end{figure}\clearpage}
\newcommand{\umllandscape}[4]{\clearpage\begin{landscape}\thispagestyle{plain}\begin{figure}[p]\centering\includegraphics[width=244mm,height=145mm,keepaspectratio]{\detokenize{../deliverables/UML_color/#1}}\caption{#2}\label{#3}\par\vspace{2mm}{\small #4}\end{figure}\end{landscape}\clearpage}
\hypersetup{pdftitle={酒店预订系统：软件设计与建模},pdfauthor={王若璠},pdfsubject={软件工程实验二}}
\begin{document}
\begin{titlepage}
\thispagestyle{empty}
\begin{center}
\vspace*{20mm}
{\Large 软件工程课程实验}\par
\vspace{18mm}
{\huge\bfseries 实验二：软件设计与建模}\par
\vspace{10mm}
{\LARGE 酒店预订系统}\par
\vspace{4mm}
{\large 从需求规则到可操作的界面}\par
\vspace{26mm}
\begin{tabular}{rl}
姓名：&王若璠\\[3mm]
学号：&241220103\\[3mm]
需求基线：&实验一酒店预订系统需求规格说明\\[3mm]
日期：&2026 年 11 月 10 日
\end{tabular}
\end{center}
\vfill
\end{titlepage}
\pagenumbering{roman}
\tableofcontents
\clearpage
\pagenumbering{arabic}

\section{项目背景与推进过程}
\subsection{这次要接着做什么}
我选的题目是校内的酒店预订系统，面向校内人员和校外住客，覆盖房源查询、在线预订和前台接待三块业务。实验一交上去的是一份二十八条的需求说明，其中十九条标了已确认，剩下九条还没进行，主要是价格筛选的口径、并发时的细节裁决和一部分接待流程。实验二的任务是把这些文字落到模块、对象、流程和界面上，让后面的实现有东西可照着写。

那九条没确认的需求，我是另开了一份工作资料，把这类没定论的地方编成 A01 到 A16 十六个候选方案，画图和写界面时直接引用编号，附录里有对照表。这样一眼就能看出哪条规则是甲方原话，哪条是我自己先垫上的假设。（上课的时候老师也说过，后续实现和这次设计对不上是正常的，就没必要在报告里把假设说成结论）

\subsection{我按什么顺序做的}
最主要的顺序是四步：先理模块、状态和数据关系，再画 UML，然后做界面原型，最后回头查。（虽然很简洁但是我中间还是返工了两次，第二次返工还是被一个输入框逼出来的。

第一步我是先想清楚订单到底有几个状态，房间到底按什么单位记。一开始我把有没有房当成一个现成的数字，后来发现不太对，同一个房型，只有当天的容量减掉当天的有效占用，才算得出某个晚上还能卖几间；后来又发现“这间房空着”和“这间房现在能住人”根本不是一回事，只能再拆出一层具体房间。我在这一步反复改了三四轮（

第二步按 PlantUML 的语法画图，一共九张，覆盖六种类型，见第\ref{sec:uml}节。第三步做界面原型，用 HTML/CSS/JavaScript 写成浏览器里能点的一整套页面，不接后端，数据就存在当前会话里。第四步用 Playwright 驱动 Chromium 把主要路径跑一遍，把截图和检查结果存下来。

第一版原型做出来的时候其实还算完整，十八组页面、七十二张状态截图，页面能切、弹窗能开、桌面和手机视口不横向溢出，自动检查全绿。但我自己点进去发现很别扭，主要是ui做得太差劲了，逻辑有问题，美观也很不好，比如说绝大多数按钮只会弹一句说明，根本不知道该怎么一步一步操作，反正就是非常混乱，然后我把这些问题全都列了出来，让codex整体重做。

第二版我想的就是必须能让使用者清晰明白地操作。顾客端从查房进入，登录以后回到原来那张预订表单，填过的东西还在；前台和管理员各自有固定工作台，不再和顾客页面混在一起；预订、接单、订单查询共用同一份演示数据，前台确认了订单，顾客那边刷新就能看到已确认等等等等。

然后第二次返工主要就是修补一些小的漏洞。比如我发现入住人数只能从一到六里选，六个人以上的团就没办法操作了，然后就让codex把人数改成整数输入等等。第\ref{sec:ui}节细说。

\subsection{三类使用者和功能边界}
顾客关心的是能不能找到合适的房、付完钱以后订单走到哪一步；前台每天要接单、接待、盯房态；管理员主要看资料变更和平台整体情况。这三类人的操作频率和关注点差得比较多，所以界面入口和权限范围我也就分开设计了。

\begin{table}[H]\centering\small
\caption{三类使用者的主要任务和权限边界}
\begin{tabularx}{\linewidth}{p{19mm}YY}\toprule
角色&主要功能&不能越过的线\\\midrule
顾客&注册登录；按条件组合查房；模拟付款；查看、取消和修改自己的预订&只能动本人的订单；付了钱也只是“待接单”，还得等酒店确认\\
酒店工作人员&本酒店接单；线下订房；入住退房与清洁；提交资料和房量价格变更；无法接待时的处置&只处理自己所属的酒店；资料和常规房量变更要经管理员审核\\
平台管理员&审核变更申请；设置候选的超时费率；跨酒店查订单、看统计和操作日志&能全局查，不等于能替顾客取消订单或者替前台接单\\\bottomrule
\end{tabularx}\end{table}

（这个原型里的付款和退款都是模拟的，没有接任何真实支付。前台给客人打电话确认、向公安备案这些事在系统之外完成，会员、积分和第三方平台同步也没有放进这次范围。界面里的“栖校”是我用来演示编的名，梧桐院、青禾馆、书庭居三家酒店也都是虚构的，照片也是让codex生成来示意房间长什么样。）

\section{架构和关键规则}
\subsection{为什么是模块化单体}
我在架构上选的是模块化单体：就是界面调应用服务，应用服务组织领域规则和数据访问，业务记录落在关系数据库里。把模块职责分开，但初期不拆成独立部署的服务。

理由也很简单，一次订单确认同时要动订单状态、好几个入住夜的房量、还有一条操作记录；客人取消的时候还要连带一笔退款。这些变化要么一起成功，要么一起失败，放在一个数据库事务里最简单，拆成几个服务反倒要处理跨服务的数据不一致。（当然这个假设也有个前提，就是付款是模拟的，账本和订单在同一个库里。以后真接第三方支付，外部交易成功而本地提交失败这种情况就得另外设计补偿了，不能拿现在这个结论套了。）

\begin{table}[H]\centering\small
\caption{九个模块的分工}
\begin{tabularx}{\linewidth}{p{9mm}p{30mm}Y}\toprule
编号&模块&负责什么，跟谁打交道\\\midrule
M01&身份与授权&认人，查角色、订单归属和酒店范围；所有入口都要过它\\
M02&酒店资料与审核&保存已生效的资料和待审的变更；审核房量变更时回头调库存校验\\
M03&库存与报价&逐夜查可用量、出报价、占房和放房，顺便管停售\\
M04&订单协调&组织创建、确认、取消和修改，需要时调库存、款项和日志\\
M05&模拟款项&记付款和退款，保证同一笔钱不会被记两次\\
M06&前台接待&线下订房、核验身份、分房、入住退房和清洁\\
M07&自动任务与应急&发起超时取消、未到店结束和故障处置，复用订单那边的规则\\
M08&查询与统计&在授权范围内查订单，算口径说得清的指标\\
M09&操作记录&记下谁在什么时候对什么做了什么、结果如何\\\bottomrule
\end{tabularx}\end{table}

订单协调模块是整个流程的组织者，但它不管理细节。它发起确认订单，库存模块来判断某一个晚上房够不够，款项模块就负责退款记录。后台的定时任务也一样调这些能力，不然人工接单和自动取消很容易各维护一套判断，两边口径慢慢就分叉了。

\subsection{“这间房有空”和“这间房现在能住”}
这一块是我改得最多的地方。每日库存是说某个住宿日期还能确认几间房，具体的房间状态是说现在能不能把这间房交给客人。今天刚退房还没打扫的房间，不影响它明天以后能不能卖；反过来，打扫干净的空房，也不代表未来每个晚上都没人订。

所以我分了三层：\pathtext{InventoryDay} 按房型和入住夜记容量、停售和版本，\pathtext{Allocation} 记订单产生的有效占用，\pathtext{PhysicalRoom} 记房号和当前的接待状态。入住时再用一条记录把订单、具体房间和入住人串起来。前台办入住只做接待分配，不会再把已经确认占用的房量扣第二遍。

订单和款项也分开记。取消不一定退款，比如已经确认的单子在入住当天十八点之后取消就不再退了，所以要退多少必须看原来的实付、退款原因和原账户，不能凭一个已取消判。

\begin{table}[H]\centering\small
\caption{核心对象以及为什么留这些字段}
\begin{tabularx}{\linewidth}{>{\raggedright\arraybackslash}p{43mm}Y}\toprule
对象&留它的理由\\\midrule
Account、Hotel、RoomType&装用户身份、酒店归属和房型容量，权限和查询都要用\\
Booking、BookingItem&订单存渠道、联系人、日期和状态；明细存房型和数量，不必把一单永远限死成一间\\
NightPrice&把下单那一刻逐夜的成交价留住，以后改价也不改历史订单\\
InventoryDay、Allocation&每晚核容量和占用，放房只放一次\\
Payment、Refund&收款和退款分开记，退款挂在原付款上，总数不超过实付\\
PhysicalRoom、Stay、GuestRegistration&管房号、实际入住和全部入住人，和预订库存互不干扰\\
ChangeRequest、ModificationRequest&一个管审核版本，一个管改单过程中原单和新单的关系\\
OperationLog&追溯关键操作，但不往日志里塞密码和完整证件号\\\bottomrule
\end{tabularx}\end{table}

我在类图里留了一单多明细的结构，现在的界面先只演示同一房型订多间。混合房型组单没有因为图上画了明细就顺手算完成。另外像“房型加日期”的唯一约束、外键、事务、请求键、金额用整数分存这些，都是留给数据库实现阶段落实的事。

\subsection{按夜核对，而不是只看入住那天}
住宿区间包含入住日、不包含退房日。周五住、周日走，占的是周五和周六两晚，周日晚不能算进去。对同一房型，设第 $d$ 晚的可用容量是 $C_d$，已有有效占用 $A_d$，要订 $q$ 间，就得在整个区间里满足
\[
q\leq\min_{d\in[\text{入住日},\text{退房日})}(C_d-A_d),
\]
而且每一晚都不在停售范围里。只要有一晚不够，这段住宿就不能显示为可整段预订。实现时还得保证多晚的占用一起提交，不能第一晚扣成功、第二晚失败，最后留下半笔订单。

待接单订单不占房，这是一开始的需求里已经了确认的。意思是说“查询显示有房”和“付款成功”都不能直接承诺房量，确认的时候必须重新读一遍订单、付款时间、全部入住夜和同时段的其他订单，再决定是确认还是退款。

\subsection{取消、修改和无法接待}
\begin{table}[H]\centering\small
\caption{取消时的退款与房量处理}
\begin{tabularx}{\linewidth}{p{52mm}p{29mm}Y}\toprule
情形&退款&房量和后续\\\midrule
待接单期间顾客自己取消&全额&本来就没占房，库存不动\\
付了钱满三十分钟还没被接单&全额&系统取消，之后不能再按待接单确认\\
确认的时候已经没房了&全额&不产生有效占用\\
已确认、没入住，入住日十八点前取消&全额&释放这一单还生效的占用\\
已确认、没入住，入住日十八点整及以后取消&不退款&占用照样释放\\
酒店原因无法接待，已确认且未入住&全额&释放预订占用，但故障停售保留\\\bottomrule
\end{tabularx}\end{table}

修改预订我是按先取消原单、然后重新付款、最后等确认来处理。这么做有个后果就是原单取消掉以后，新单可能因为没房下不了，原房也就回不来了。所以界面在动手之前得把原单能退多少、新单要付多少分别列清楚，把风险弹出来让客人知道（类似弹窗警告那种）；后台需要存下来新旧订单的关联和修改进度，方便查结果和安全重试。

无法接待又不一样，出故障以后第一步是先把受影响的房源停止销售，别让新的订单继续进来，第二步才是逐单核对。就算退款把占用释放了，故障还是没有消失，所以“释放库存占用”和“恢复可售”必须是两个动作，不能一退款就自动重新开卖。已经入住的、只受局部影响的订单让人工处理。

\subsection{留给实现阶段的三个保护}

1. 资源授权。把按钮藏起来只是界面好看，挡不住越权，正式系统每个请求都得查这张订单是不是这个顾客的，这家酒店是不是这个工作人员的。管理员能看全局，也不等于能替别人做写操作。

2. 并发下的重读和一次提交。候选方案是在同一家酒店内用一把统一的库存协调锁，接单、取消、线下订房、审核和应急处置都按同样的顺序拿锁，再按订单号读相关记录。这样同酒店的写操作会排队，并发度是降下来了，但初期更容易验证正确性。这套东西有没有用，还得到实验三的时候里拿真正的同时请求和失败回滚试。

3. 防重复处理。用请求键识别同一次操作，已经成功的请求原样返回结果，同一个键配不同内容要拒绝。两个不同的预订，不会因为顾客和日期相同就被合并成一条。退款、释放占用、后台补扫和修改重试都要遵守这一条。

\section{UML 建模及各图的说明}\label{sec:uml}
我这次一共画了九张图，覆盖六种类型：结构性图里选了类图和组件图，行为性图里选用例图、状态图和活动图，交互性图选了时序图。每张图都留了 \pathtext{.puml} 源码和导出的 PNG。任务书要求三类各至少一张、类图必画。（我是先想清楚每张图要回答哪个问题，再去挑图种）

\begin{table}[H]\centering\small
\caption{九张图各自回答什么问题}
\begin{tabularx}{\linewidth}{p{14mm}p{25mm}Y}\toprule
编号&图类型&回答的问题\\\midrule
D01&用例图&谁可以做什么，范围划到哪儿\\
D02&组件图&九个模块怎么配合，数据从哪儿过\\
D03&类图&订单、房型、逐夜价格、房量和款项怎么连\\
D04&活动图&改单的每一步顺序，失败以后还留下什么\\
D05&状态图&订单允许哪些转换，哪些状态回不去\\
D06&时序图&接单那一瞬间怎么重查超时、竞争和库存\\
D07&时序图&取消时怎么同时处理状态、占用和退款\\
D08&活动图&核验、入住、退房、清洁怎么接上\\
D09&活动图&出故障以后怎么停售、退款、留人工待办\\\bottomrule
\end{tabularx}\end{table}

\subsection{谁在什么范围里做什么}
用例图（图\ref{fig:d01}）把顾客、酒店工作人员和平台管理员分成三块。顾客注册是单独一个用例，不能从注册入口里选个管理员身份出来。工作人员那几个用例（接单、维护资料、应急处置）都圈着本酒店这个范围。（管理员主要连到审核和全局查询，我不想把他画成一个不受限制的万能角色。）

组件图（图\ref{fig:d02}）往下走一层，说明这些操作落到哪个模块。订单服务协调库存和款项，资料审核在生效前要回头查库存，自动任务调的是同一套订单规则。鉴权是所有入口都要有的一条约束，但我在图里没有给每个组件都拉一条鉴权线，那样连线会把主体关系盖掉，只在说明里提了一下。

类图（图\ref{fig:d03}）只表达业务记录的关联，有三组关系最重要：订单带明细和逐夜成交价，明细产生每夜的占用，付款关联退款。具体房间和实际入住单独挂一条线。这么连的好处是能从一笔订单一路追下去，订了什么、占了哪几晚、付了多少钱、后来退没退。

\subsection{订单能走到哪些状态}
状态图（图\ref{fig:d05}）第一个切在“待接单”和“已确认”中间。在线付款成功只进待接单，前台确认了才真正占房。线下订房在候选方案里是收款和确认占用一起做，所以可以直接进已确认。入住后转已入住，正常退房变已退房；取消、退房、未到店结束这三个终态，都不会因为旧页面又提交一次就变回有效预订。

“未入住结束”和“已退房”我一开始想合成一个“结束”，后来发现不行，没来住过的人从来没进过房间，为什么给他生成一次清洁任务？按需求还没到店的订单还保留原住宿期，到原定退房日十二点才结束，不能因为第一晚没来就把后面几晚提前卖掉。这个后台流程原型里还没做，状态图留着它是给实现阶段看的。

修改活动图（图\ref{fig:d04}）先把输入、原单退款和新报价核对一遍，再让顾客确认风险。原单取消和新单付款是两个阶段，中间不能跨着顾客一直持锁。如果原单已经取消成功，新单付款或者接单失败，结果页要把两笔订单各自的结局都展示出来。

\subsection{接单那一瞬间，以及取消}
接单时序图（图\ref{fig:d06}）把正常确认、已经超时、没房、版本冲突和重复请求放在同一条流程里。它解决的是前台看到订单的时候有房，不代表点下确认按钮的时候还有房。所以订单服务必须在受控的事务里重新判断一次，然后把占用和状态一起写进去。如果这一单确认下去会让别的竞争订单没法履行，还得准确找出受影响的那些待接单订单，全额退款。

取消时序图（图\ref{fig:d07}）分两支。待接单没有占用，取消就是改状态加全额退款；已确认未入住要按入住日十八点这个界线算退款，同时释放这一单还在生效的占用。图里把相关写入放在同一个数据库事务里，就是为了不出现订单已取消但房没释放或者退款记了账但取消失败这种情况。

两张图都画了重复请求直接返回已有结果的分支。后台重试和用户连点应该复用已经提交过的结果，不再退一次款、再占一次房。

\subsection{前台的一天，和出事的时候}
接待活动图（图\ref{fig:d08}）把核验身份、登记全部入住人、挑具体房间、更新订单串起来。退房以后房间先进待清洁，打扫完才回到可安排。续住、提前退房和延迟费用这几件事还有没定的条件，所以图里标的是转人工核对或者走候选收费方案，界面没有替这些未知规则拍板。

无法接待活动图（图\ref{fig:d09}）先封住影响范围，再一单一单核对。已经取消或者退过款的订单不能重复处理，已入住和只受局部影响的留成待办。这里的循环条件用的是还有没有未核对的订单，不是所有人工事项都处理完了，不然的话只要有一条待人工的记录挂着，整个流程就永远退不出去，图的意思也就废了。

\umlportrait{D01_用例图_角色权限.png}{用例图：三类角色和各自的权限范围}{fig:d01}{对应需求 REQ-002、003、004、011、015、023、025。}
\umllandscape{D02_组件图_整体架构.png}{组件图：模块化单体的主要协作关系}{fig:d02}{画的是目标架构；当前浏览器原型还没接关系数据库，也没有服务端鉴权。}
\umllandscape{D03_类图_核心领域.png}{类图：核心业务对象及其关联}{fig:d03}{多房型明细是候选结构；界面目前只做单房型多间。这张图不代表数据库已经建好。}
\umlportrait{D05_状态图_订单生命周期.png}{状态图：订单的生命周期}{fig:d05}{付款、退款、房间状态各自保存，不能只靠订单状态一个字段推出全部结论。}
\umlportrait{D04_活动图_修改预订.png}{活动图：修改预订的步骤和失败分支}{fig:d04}{现在的修改界面支持改日期和房型；图里改数量的部分属于后续完整设计。}
\umlportrait{D06_时序图_接单竞争.png}{时序图：前台接单、竞争核验与失败处理}{fig:d06}{图说明的是调用顺序和目标事务边界；锁、请求键和回滚还要服务端来实现。}
\umlportrait{D07_时序图_取消退款.png}{时序图：顾客取消与模拟退款}{fig:d07}{退款截止时间按入住日算。酒店责任的退款单独用一个原因，不套顾客自愿取消的规则。}
\umlportrait{D08_活动图_接待与清洁.png}{活动图：入住、退房与清洁}{fig:d08}{A09、A10、A16 涉及候选收费、续住和提前退房，还没到能当已确认规则用的时候。}
\umlportrait{D09_活动图_酒店无法接待.png}{活动图：无法接待时的停售与订单处理}{fig:d09}{应急停售和恢复可售分开做；退款退完了不代表故障已经排除。}

\section{界面设计与操作流程}\label{sec:ui}
界面是按任务来组织的。顾客端用顶部导航，前台和管理端用固定侧栏，三类人拿到的导航不一样。图版用的是真实浏览器截图，有几张只截了主内容或弹窗，完整的 PNG 我放在 UI 界面目录里了。

\subsection{首页、查房和人数}
首页第一眼就是查房表单：校区、入住日期、退房日期、人数。下面一排住宿卡片，卡片上给的是整段住宿的总价和“查看”入口。介绍性的照片用来辨认环境，但页面的主角始终是那张表单。
\ui[99mm]{home.png}{顾客首页：查房表单是唯一的起点}{fig:home}

人数这一项是数字输入，允许一到三十，可以点上下箭头调。空值、负数、小数、指数写法和越界都会拦下来。（这个30只是我用来演示的上限，不代表甲方说过团体最多三十人，这个改起来就是改一个常量。）

多人查询按同一房型需要的间数算。设总人数 $n$、每间容量 $c$，最低间数是 $\lceil n/c\rceil$。比如七个人选每间住三人的房型，就得三间；演示单价每间每晚六百八，住两晚是 $3\times2\times680=4080$ 元。预算筛选也得用这个总价，不能在结果页上摆一间房的低价，到了付款页才乘以三。

\clearpage
\ui[92mm]{results.png}{查询结果：日期、人数、房型和全部房间总价一起筛}{fig:results}

只要有一晚凑不齐需要的间数，这段住宿就没有可订结果，不会给你一个假装能订的选项。空结果页留了“清除筛选”的入口，不把用户堵住。人数放开以后，原型仍然只处理同房型多间，混合房型自动搭配还没做。

\subsection{详情、预订和登录回跳}
详情页把酒店介绍、设施、房型容量、住宿日期和费用堆在一页里。每个房型旁边直接写清楚这次人数至少要几间，房量不够就不让往下走。页面也保留了改查询条件的入口，不用退出去重新从首页找。
\ui[74mm]{detail.png}{酒店详情里的房型选择与行程摘要}{fig:detail}

\clearpage
预订页填联系人、电话、房间数，右边列日期、人数、逐夜价格和合计。房间数默认就是刚才算出来的最低间数，可以加，但不能减到装不下人，也不能超过各晚的可订量。预订人和实际入住人允许不同，全部入住人的信息到店时让前台核验。
\ui[82mm]{checkout.png}{预订信息、取消规则和费用核对}{fig:checkout}

\ui[67mm]{login.png}{登录页：回到查房，也留了演示账号}{fig:login}
登录和注册不会把整套顾客导航换掉。没登录的人从预订表单跳去登录，成功以后回到原来那一页，联系人、电话、间数和规则勾没勾都还在。原型支持演示账号和临时注册，不存密码，也不意味着服务端认证已经做好了。

\clearpage
付款按钮旁边写着“付款后还需酒店确认”。付完钱的页面标题是“等待酒店确认”，就是不想让用户把“付款成功”读成“订房成功”。页面上顺带解释了没房、三十分钟没接单、待确认期间自己取消这三种情况钱怎么退。
\ui[105mm]{pending.png}{模拟付款后的待确认页面}{fig:pending}

\subsection{订单查询、取消和修改}
“我的订单”分全部、待确认、待入住和已取消四类。先看清列表，点进详情才出现可用操作，终态订单不再挂一个点了也没用的取消按钮。多人订单的详情里留着预订人数和房间数，方便和自己填过的对一遍。
\ui[48mm]{orders.png}{我的订单：分类、状态和进入详情的入口}{fig:orders}

\clearpage
取消之前弹一次确认框，把原实付和这次能退多少摆在一起。图\ref{fig:cancel}这张是我专门造的一个已过入住日十八点的已确认订单，用来看零退款那一支，不是真有人住过。用户可以留着订单，也可以确认取消，操作完页面直接给结果。
\ui[59mm]{cancel.png}{截止后取消：先亮出零退款，再由顾客决定}{fig:cancel}

修改要核对两次信息：先填新日期和新房型，再看原单退多少、新单付多少、原房不保留的风险。选的房型装不下原预订人数就不许提交。改完的订单留一条回到原单的链接，原单和新单各自的付款、退款和状态分开显示。
\ui[57mm]{modify.png}{修改时的费用与风险确认}{fig:modify}

\clearpage
\ui[83mm]{modify_failed.png}{新单没成时，仍然能追到那张已取消的原单}{fig:modify-failed}
图\ref{fig:modify-failed}用的是造出来的失败数据，想说明的是原单取消成功和新预订成功就是两码事。原单不会因为新单失败又活过来，这个在确认之前要给客户展示清楚，结果页也得讲清楚，不能丢一句操作失败让人猜。

\subsection{前台接单和线下订房}
前台工作台从待接单开始，上面显示待处理、已确认、在住和可安排房间的数量。每条待处理记录都带着联系人、房型、住宿日期、实付金额和状态。点确认时会重查一次房量；如果有更早付款、也还能履行的竞争订单，先提示处理那一单。
\ui[79mm]{staff.png}{前台接单：固定导航与待处理订单}{fig:staff}

\clearpage
线下订房单独一个入口。到店的人才不用先注册平台账号，工作人员填联系人、房型和日期，逐夜查完房量，模拟收款，顺手确认。符合条件的学生要现场核验校园卡，八折是按对应日期的公开价算的，不能在网络优惠价上再打一次，也不能给已经在网上订好的单子事后补线下优惠。
\ui[80mm]{offline.png}{线下订房：联系人、日期、间数和校园卡核验}{fig:offline}

\subsection{入住、退房与清洁}
入住页面只列能办的订单和在住的订单。办理时先核验身份；如果是别人代订，前台要在系统外联系预订人，再登记实际住客、分配能接待的房间。原型不收真实证件号，正式实现还得把需求里那几个登记字段和保护措施补上。
\ui[66mm]{checkin.png}{入住登记：核验、实际住客与房号}{fig:checkin}

\clearpage
正常退房以后，那间房转成待清洁。房态页把可安排、已入住、待清洁和故障停用分开列，清洁按钮只出现在待清洁的房间上。前台点一下确认打扫完，这间房才重新能被安排给下一位。
\ui[72mm]{rooms.png}{房态与清洁：接待条件单独管}{fig:rooms}

续住的入口会先查一下后面的房量；提前退房和延迟费用没有自动套用还没确认的规则。图\ref{fig:d08}画的是完整候选流程，能点的只是其中一截，这两件事得分清。

\subsection{资料、房量价格和审核}
资料维护可以改酒店介绍和设施说明，提交后变成一条审核申请。审核通过之前，顾客那边看到的还是原来已经生效的内容，待审的东西不能直接发出去。照片上传和文件存储没做，页面用的是本地示意素材。
\ui[62mm]{catalog.png}{酒店资料维护与申请记录}{fig:catalog}

\clearpage
房量价格页按入住夜列出容量、有效占用、可售量、价格和停售状态。工作人员提交的是变更申请，管理员审核生效时得再查一次最新占用。举个例子：提交的时候想把某晚减到一间，审核前已经有两间被确认掉了，那就不能按旧的条件通过。
\ui[84mm]{inventory.png}{按日期提交房量、价格和停售变更}{fig:inventory}

最后就是审核页，显示申请内容和状态，管理员可以通过，也可以填了原因驳回。候选的延迟退房费率单独设了一个值，但原型不会拿它自动去算那些还没定的场景。常规停售如果会影响已确认订单，会被拦下来；真出了故障的话就往无法接待的流程走。
\ui[75mm]{reviews.png}{管理员审核与候选费率设置}{fig:reviews}

\clearpage
\subsection{无法接待、全局查询和操作记录}
无法接待页面先填受影响的房型、日期和原因，再预览一遍会牵扯到哪些订单。确认以后先停售，未入住的订单按酒店责任退款，在住的列进人工处理。原型现在是把所选房型整个停售来演示的，局部房间影响和恢复可售的条件还得继续做。
\ui[78mm]{incident.png}{无法接待：影响范围和处置入口}{fig:incident}

管理员的全局订单页能按酒店、状态、日期类型和关键词筛，详情用来看处理过程，没有加“代客取消”这种还没确认的按钮。日期现在给的是入住和退房，候选设计里下单、付款日期的筛选留到后面实现。
\ui[75mm]{global_orders.png}{跨酒店订单查询}{fig:global-orders}

\clearpage
经营概览显示本次会话的订单量、模拟收款、退款和净收款。订单量把取消单也算进去，修改产生的新旧订单分别计数；净收款就是收减退，所以没有把它叫成“收入”。缺了完整的房晚历史，入住率那个百分比算出来也没办法复核，所以就直接不放了。
\ui[86mm]{overview.png}{经营概览：指标口径与当前房间接待情况}{fig:overview}

操作记录列出时间、操作者、动作、对象和结果，用来把接单、清洁、审核这些动作连回到具体的业务对象上。现在的日志来自本地演示；以后服务端得保证成功的业务日志和业务提交是同一笔，失败的安全事件另外记，访问范围也要限住。
\ui[49mm]{logs.png}{关键操作记录}{fig:logs}

\clearpage
\section{检查结果与没做到的部分}
\subsection{我实际跑了哪些检查}
我的检查脚本是 Playwright 驱动本地 Chromium 跑的，脚本本身是让codex写的，跑完的记录和截图都留着。普通交互记录十五组，人数改动另外五组。（这里说的“组”里含好几步操作和断言，不是二十条后端验收用例。）
\begin{table}[H]\centering\small
\caption{有实际记录的检查及结论}
\begin{tabularx}{\linewidth}{p{34mm}YY}\toprule
查了什么&实际看到什么&能说明什么\\\midrule
登录与预订&导航没丢；登录回来表单还在；付款后进待接单&顾客这条主线能连着操作\\
跨角色接单&前台确认同一单后，顾客端显示已确认&同一会话里状态是联动的\\
取消与修改&退款前要确认；终态不显示取消；原单取消并挂上新单&不同操作的后果在界面上分得开\\
入住退房清洁&房间从可安排到在住、到待清洁，再回到可安排&接待和清洁能演一遍\\
审核与故障&容量冲突的申请生效不了；未入住退款，在住留人工&这几条前端分支确实跑过\\
人数与多间房&非法输入拦住；七人三间报价一致；一晚不够就不给订&人数改动确实传到了查询和结算\\
页面与素材&十九条路由在桌面和手机视口都没有整页横向溢出；图片都加载出来了&基础布局和素材引用正常\\\bottomrule
\end{tabularx}\end{table}

UML 的 PNG 里嵌着源码，我把它和对应的 \pathtext{.puml} 逐个对过，确认真的是同一版。二十六张 UI 截图保留了路由、构造数据的说明和文件摘要。完整的证据我放在 checks 目录了，附录里列了入口。

\subsection{这些检查证明不了什么}
第一版所有渲染检查都过了，我照样能指出导航断了、任务没入口。这件事本身就说明检查范围决定了结论：页面能切、弹窗能开，跟用户能不能把一单订完完全是两码事。第二版加了任务路径检查，是针对这个漏洞补的，但也远没覆盖全部业务。

设计上还挂着的边界包括混合房型组单、照片上传、完整证件登记、续住结算、提前退房、部分在住的故障和统计口径。有的在模型里留了位置，有的还等业务确认。这些之后还得一项一项定。

\section{AI 辅助过程}
\subsection{大致过程}
我用 Codex 做资料整理、方案生成、PlantUML 起草、前端原型修改和检查脚本编写；真正负责渲染和编译的是 PlantUML、浏览器自动化和 XeLaTeX。下面按六次主要的使用场景来讲，一次场景里往往包含好几轮来回，完整对话以实际会话为准，这份索引不能冒充每一条原始输入输出。

我在其中提的要求主要是三类：设计要细、要分阶段给我看；第一版的导航、页面编号、任务入口和观感都不行，重做；入住人数不该只能选到六人，改成有范围的整数输入。这些反馈确实推动了修改。至于架构取舍、代码自查和自动检查，凡是工具提出或执行的，下面也照实记在工具那一栏。

\subsection{第一次：把任务拆开}
第一次是让工具梳理实验二的任务，以及它和实验一是什么关系。输入是任务书、两份往年示例和已有的需求，输出是这次要交哪些图种、哪些模块界面、过程记录和提交材料。它的用处是帮我把任务拆成能排期的几块，顺带发现任务书里的截止日期是旧的。局限也明显：分析结论还得回原文核，工具列出来的东西很多并不都是课程的硬性要求。

\subsection{第二次：定架构和业务边界}
第二次让它出架构和业务边界的候选，产出了九个模块、十六项假设、需求映射和五十二条验证计划。我觉得比较有价值的一点，是把订单、款项、逐夜库存和具体房间分开——这也是后面所有图和界面的地基。风险在于，工具很爱把没说明的细节补成一条完整规则，看着特别顺。所以资料里一直挂着“候选”“待确认”，没有把生成的方案倒过来当成甲方的答复写进需求。

\subsection{第三次：把设计变成图}
第三次是把设计换成 PlantUML 代码。工具出了六种九张图，中途还处理了标签换行引起的语法错误。后面我又让它改了几处：顾客注册从共同登录里拆出来单独成用例；接单、取消、审核、应急的库存协调范围统一到酒店一级；应急流程里那个循环的含义也重画了，原来看不出它要退出。

用代码画图的好处很直接：改名、对照、重新导出都方便。但能编译只说明语法和渲染没问题，关系对不对它管不了。拿修改预订举例，如果图里把“新单失败”画成原单恢复，图片再整齐也是错的，因为原房不保留是需求里写死的。所以图还是得回到状态、占用和款项的实际含义上去解释。

\subsection{第四次：第一版界面为什么不能用}
第四次让它出第一版界面，十八组页面、七十二张状态截图。工具自己检查了页面切换、弹窗和布局，一条都没检查能不能订完一单。我看见原型以后提的问题很具体：进登录页左边只剩一个选项，顶部的 U10、U02 不该出现在产品里，页面乱，不知道从哪开始。

这次暴露的是自动生成的一个很大的问题：它能快速堆出很多看着很完整的页面，但组织方式是把设计成果展示出来，不是让顾客干成一件事。页面编号对开发和需求映射是好事，对住客是噪音；按钮能弹个窗，也不等于业务流程走得通。

\subsection{第五次：按意见重做}
第五次是照着上面那些负面反馈重做。顾客端改成查房、详情、预订、订单一条链，前台和管理员分开，页面编号挪进独立的设计资料，按钮真的去改本地演示数据。改完又补了登录回跳、前台接单、取消修改、入住清洁、审核这些路径的检查。

这一轮里，我说的是“原版不好用，重做”，工具负责改和复查。它自查时也确实揪出过几个问题：登录返回时字段丢了、最低起价算错、改酒店资料影响到了别的酒店、还有一个变量在订房时用错了地方。把这些修正记下来，比笼统写一句“AI 提高了效率”有用得多——它确实让试做变快，也确实会带进需要人盯着的问题。

\subsection{第六次：一个人数输入框牵出的连锁修改}
第六次是从我那句话开始的：入住人数怎么只能选到六人。表面上是换个控件，工具一查才发现，查询那里拿总人数去比单间容量，超过一间的人数整家酒店就被筛没了。光把下拉框放宽到三十，七个人还是订不了三间。

后面的改法是把人数做成一到三十的整数输入，间数按房型容量向上取整，查询、详情、结算统一用全部房间的总价，订单多存一个预订人数。新加的检查覆盖非法值、上下界、七人三间、登录回跳后数量还在、以及某一晚库存不够的情况。

这里给我的印象还是比较深的，界面字段不是孤立的。一个人数能牵到容量、间数、价格和结果的保存。工具在找关联位置上是帮得上忙的，但前提是我自己先把“人数”到底指单间人数还是整笔订单人数想清楚，否则改得越快，错的口径传播得越远。

\subsection{用完之后的一点体会}

第一，大模型最擅长的还是把重复的、格式化的活干掉：拆任务、出候选方案、生成能改的图和能点的原型、写检查脚本。这些东西它做得比我快，而且我随时可以改。

第二，它报的完成和通过必须连着检查范围一起看。第一版七十二张截图一张不缺，如果一次导航断路了那全完蛋了，不能光是数字好看。

第三，反馈说得越具体，修改越有效。“再完善一点”这种话，工具只能给我换一批相似的页面，如果给它说清楚“哪个入口消失了”“哪个字段挡了真实任务”，会高效很多很多。

\clearpage
\section{阶段结果与下一步}
这一阶段拿到了九个模块的职责划分、六种九张 UML 图、三类角色的可交互原型和二十六张界面图。设计上重点处理了几件容易出错的事：付款和确认分开、房量按夜核、取消退款的分界、修改后原房不保留、接待与清洁的衔接、故障时停售和恢复可售分开。第二版原型又照着具体反馈改了导航和人数查询，界面和业务规则的联系比第一版清楚不少。

下一步先把订单、库存和模拟款项落到数据库，把事务、唯一约束和资源授权补齐；然后做后台的超时取消和未到店结束，再验一遍重启后的恢复。并发接单、同一请求重复提交、逐夜写入中途失败、越权访问这几件事，应该排进实验三最先验证的位置。还悬着的业务条件继续找来源、定方案。

这一阶段最重要的吧，我感觉是几个不变量，待接单不占房、房量按夜核、原单取消后原房不保留、停售和恢复可售分开。这几条只要在实现时不破，别的细节怎么改都还能兜回来。

如果让我重新做一遍，我会先把“一次任务走得通”写成检查项，再去铺页面。

剩下要做的，前面那一节已经列了清单。那部分得留到实验三，边写边试。

\appendix
\clearpage
\section{图中候选设计条件索引}
图上的编号用来说明这条规则还没完全确认。下表是 A01 到 A16 的概览，详细定义在设计工作资料里。
\small
\begin{longtable}{p{12mm}p{143mm}}
\toprule 编号&候选方案或保留条件\\\midrule\endfirsthead
\toprule 编号&候选方案或保留条件\\\midrule\endhead
A01&查询价格按整段住宿总价算，人数按总容量检查；具体口径仍是候选。\\
A02&数据模型允许同酒店、同日期出现多个房型明细；界面先做单房型多间。\\
A03&金额以整数分保存，逐夜应用优惠并舍入，再按间数求和。\\
A04&竞争范围内按付款时间和订单号排先后；并列和跨区间怎么算还要核定。\\
A05&业务时区候选为 Asia/Shanghai，正式裁决用服务端受控时间；到期和状态变化在锁内重读。\\
A06&未付款阶段用结算草稿表示，付款成功后建待接单订单；草稿保存多久还没定。\\
A07&线下顾客可以没有账号，只存联系人快照；收款和确认占用算一次业务提交。\\
A08&修改时保存原单、新单和步骤；原单取消成功后，新单失败不恢复原单。\\
A09&候选延迟费用从原退房日十二点起按不足一小时向上取整，用入住时的费率版本；细节未确认。\\
A10&续住先查后续房量，计价方式和原单怎么更新仍待决。\\
A11&本酒店前台可以发起应急停售，属候选权限；已入住和部分受影响的订单留人工处理。\\
A12&登录名只是演示标识，系统订单号唯一；历史引用保留，数据保存期限待定。\\
A13&订单查询要区分下单、付款、入住和退房四类日期；界面现在只演示后两个。\\
A14&订单量、入住率和模拟净收款要分开看；历史房晚、收入确认和统计期间的口径还需明确。\\
A15&审核可以驳回并给原因，改完重新提交；旧的生效资料单独留着。\\
A16&正常退房和未到店结束分开处理；提前退房不直接套用未到店或普通取消的规则。\\\bottomrule
\end{longtable}
\normalsize

\clearpage
\section{项目仓库与版本说明}
本项目继续使用实验一的 GitHub 仓库保存进度。实验二的设计说明、UML 源文件和图片、可交互原型以及检查记录，已经整理为一个阶段版本，方便实验三在此基础上继续实现，也便于之后回看本次设计。

\noindent 仓库地址：\par
\noindent\href{https://github.com/ffcchhgg/hotel-booking-requirements}{\pathtext{https://github.com/ffcchhgg/hotel-booking-requirements}}

\medskip
\noindent 设计与原型版本：\href{https://github.com/ffcchhgg/hotel-booking-requirements/commit/0070cfe8628b76083660e73c900c99d5d788afa9}{\pathtext{0070cfe8628b7}}。

这个版本保存的是实验二的设计和浏览器原型；数据库、服务端权限和后台任务仍属于后续实现内容。报告中的功能说明、设计理由和界面流程可以独立阅读，仓库用于补充查看原始材料及操作演示。

详细的文件目录、材料索引、原型运行方法和检查记录说明集中放在仓库首页的 README 中。需要核对某张图的源码，或查看某项检查具体做了什么时，可以从那里进入；报告和最终提交材料则在后续提交中依据上述设计版本整理。

\end{document}
